\documentclass[11pt,a4paper]{article}
\usepackage[margin=2.5cm]{geometry}
\usepackage{mathpazo}
\usepackage{enumitem}
\usepackage{hyperref}
\hypersetup{hidelinks,colorlinks=false}
\setlength{\parskip}{0.8ex}
\setlength{\parindent}{0pt}
\pagestyle{empty}

\begin{document}

\begin{flushright}
\normalsize
\textbf{Manuscript:} Sample-Size Phase Transition Ends the Hard-versus-Soft Clustering Debate\\
Hard-versus-Soft Clustering Debate in Causal Discovery\\[4pt]
\textbf{Corresponding Author:} Shuaidong Gao\\[2pt]
\textbf{Email:} gsd3247186514@gmail.com
\end{flushright}
\vspace{1em}

\begin{center}
{\LARGE\bfseries Suggested Reviewers}
\end{center}

\bigskip

We respectfully suggest five reviewers covering causal discovery
methodology, clustering theory, computational genomics, and
statistical rigor---the four pillars of our contribution. All
contact information is publicly available from institutional
directories and individual homepages.

\bigskip

% --- Reviewer 1 ---
\noindent\textbf{1.\quad Caroline Uhler}\\
\textit{MIT EECS \& Broad Institute of MIT and Harvard}\\[2pt]
Email: \texttt{cuhler@mit.edu}\\[2pt]
Prof.\ Uhler directs the Eric and Wendy Schmidt Center at the
Broad Institute and is a leading figure in causal discovery for
genomics. Her work spans graphical models, causal inference, and
gene regulatory network reconstruction from single-cell and bulk
transcriptomic data (e.g.,~\textit{PLoS Comp.\ Biol.}, 2026).\,
As a researcher at the intersection of causal discovery
methodology and cancer genomics, Prof.\ Uhler is uniquely
qualified to evaluate both our methodological contribution and
the biological significance of our pan-cancer findings.

\smallskip

% --- Reviewer 2 ---
\noindent\textbf{2.\quad Kun Zhang}\\
\textit{Carnegie Mellon University / MBZUAI}\\[2pt]
Email: \texttt{kunz1@andrew.cmu.edu}\\[2pt]
Prof.\ Zhang is an internationally recognized causal discovery
researcher whose work on ICA-LiNGAM, causal representation
learning, and domain adaptation provides an independent
perspective on our NOTEARS-based framework. Unlike
NOTEARS-lineage researchers, his expertise lies in functional
causal models and non-Gaussian methods, enabling an impartial
evaluation of whether cluster-aware gating generalizes beyond the
continuous-optimization paradigm. His recent work on causal
discovery from gene expression data (\textit{arXiv}:2606.00568,
2026) further confirms his relevance to our empirical setting.

\smallskip

% --- Reviewer 3 ---
\noindent\textbf{3.\quad Marina Meil{\u{a}}}\\
\textit{University of Waterloo}\\[2pt]
Email: \texttt{mmp@uwaterloo.ca}\\[2pt]
Prof.\ Meil{\u{a}} holds a Canada Research Chair in Reliable
Structure Discovery and is one of the world's foremost authorities
on clustering theory. Her foundational work on spectral
clustering, clustering consistency, and manifold-based methods
makes her the ideal reviewer for a paper whose central claim is
resolving a two-decade debate between hard and soft clustering.

\smallskip

% --- Reviewer 4 ---
\noindent\textbf{4.\quad Fabian J.\ Theis}\\
\textit{Helmholtz Munich \& Technical University of Munich}\\[2pt]
Email: \texttt{fabian.theis@helmholtz-munich.de}\\[2pt]
Prof.\ Theis is a leading computational biologist and a core
developer of the Scanpy ecosystem for single-cell genomics.
His work on single-cell data integration, dimensionality
reduction, and deep generative models for transcriptomics
bridges methodology and biomedical application. His expertise
in single-cell RNA-seq analysis makes him uniquely suited to
evaluate our PBMC replication, cross-modal phase transition
evidence, and the practical $n/d \approx 4$ decision rule
for genomics applications.

\smallskip

% --- Reviewer 5 ---
\noindent\textbf{5.\quad Jonas Peters}\\
\textit{University of Copenhagen}\\[2pt]
Email: \texttt{jonas.peters@math.ku.dk}\\[2pt]
Prof.\ Peters is a leading causal inference methodologist whose
work on Invariant Causal Prediction and distributional
robustness provides the statistical rigor needed to evaluate our
causal claims. His perspective complements the other four
reviewers by assessing whether our empirical findings satisfy the
statistical criteria for causal (rather than merely
associational) interpretation in observational transcriptomic
data.

\bigskip
\vspace{1em}

\begin{center}
\rule{0.4\textwidth}{0.4pt}
\end{center}
\vspace{0.5em}

We confirm that none of the above five individuals have known
conflicts of interest with this work. All are independent of the
authors and their institutions; none has co-authored with any
author within the past three years; and none has served as an
advisor or advisee of any author.

\smallskip

\noindent\textit{Note:}\ All institutional email addresses are
publicly verifiable through the respective university directories,
individual CVs, and Google Scholar profiles. Should the editor
require alternative candidates, we are happy to provide additional
names on request.

\end{document}
